% Version 2, 13 October 2026; supersedes version 1
\documentclass[10pt]{article}
\usepackage{amsmath,amssymb,amsthm}
\usepackage[T1]{fontenc}
\usepackage{lmodern}
\usepackage[margin=1in]{geometry}
\usepackage{hyperref}
\allowdisplaybreaks
\newtheorem{theorem}{Theorem}[section]
\newtheorem{lemma}[theorem]{Lemma}
\theoremstyle{remark}
\newtheorem{remark}[theorem]{Remark}
\newcommand{\lcm}{\operatorname{lcm}}
\newcommand{\B}{\widetilde B}
\title{An explicit upper bound for reciprocal sums of sets with large
pairwise least common multiples}
\author{Haoyu Chen\\ Independent researcher}
\date{Version 2, 13 October 2026; supersedes version 1}
\begin{document}
\maketitle
\begin{abstract}
For a positive integer $n$, let $\rho(n)$ be the maximum of
$\sum_{a\in A}1/a$ over sets $A\subseteq\{1,\ldots,n\}$ satisfying
$\lcm(a,b)>n$ for distinct $a,b\in A$. We prove
$\rho(n)\le 1.00189053+4162.91/n$ for every integer $n\ge1$.
Two smaller certificates give $1.00271168+638.004/n$ and
$1.00447967+133.853/n$. The proof assigns nonnegative weights to
multiples, using integer cells and residue classes modulo $420$, $60$
and $12$. Exact integer checks and an explicit Euler--Maclaurin
remainder certify all constraints. As an improvement of the asymptotic
leading constant, this improves the strongest bound we located in a
literature search of 12 October 2026, including version 1 of this note
and Chen's 1996 bound. The larger error coefficients require separate
finite-$n$ comparisons, which we give exactly. The conjecture
$\rho(n)\le1+o(1)$ is not proved.
\end{abstract}

\section{Statement and context}
A set $A\subseteq\{1,\ldots,n\}$ is \emph{lcm-separated} if
$\lcm(a,b)>n$ whenever $a,b\in A$ are distinct. Write
\[
 \rho(n)=\max\left\{\sum_{a\in A}\frac1a:
 A\subseteq\{1,\ldots,n\}\text{ is lcm-separated}\right\}.
\]
The singleton $\{1\}$ is allowed, so $\rho(n)\ge1$.

\begin{theorem}\label{thm:main}
For every integer $n\ge1$,
\begin{align*}
 \rho(n)&\le1.00189053+\frac{4162.91}{n}, &&(L=420),\\
 \rho(n)&\le1.00271168+\frac{638.004}{n}, &&(L=60),\\
 \rho(n)&\le1.00447967+\frac{133.853}{n}, &&(L=12).
\end{align*}
The exact coefficients $K$ of the $1/n$ terms before the upward decimal
roundings are, respectively,
\[
 \frac{2081454577}{500000}=4162.909154,\qquad
 \frac{127600737}{200000}=638.003685,\qquad
 \frac{2091444}{15625}=133.852416.
\]
\end{theorem}

We prove this theorem directly from three finite arrays of nonnegative
rational numbers. Every error term in the proof is explicit; no limit
optimization or approximation theorem for a class of functions is used.
The arrays have $39$ rows, with $24$, $12$, and $6$ columns respectively.
They are supplied as JSON files, together with the verification programs.

\paragraph{Related work.}
Problem 542 in Bloom's list \cite{EP542} includes the classical bound
$31/30$ and a separate divisor question. The conjectural strengthening
$\rho(n)\le1+o(1)$ remains distinct from those questions. The literature
records $\rho(5)=31/30$ and $\rho(11)=4699/4620$ \cite{Ch96,We}.
Chen \cite[Theorem 2, p.~104]{Ch96} proved, explicitly,
\begin{equation}\label{eq:chen}
 \rho(n)<B_C(n):=1.0170166+\frac{14.21}{n-58}\qquad(n\ge59).
\end{equation}
Weighted counting over disjoint sets of multiples already appears in
that paper (Theorem 1 and Lemma 1). The weak duality principle used here
is therefore prior methodology. Our contribution is the three specific
periodic arrays and their certified numerical costs. Chen's later
paper \cite[Theorem 1]{Ch97} treats elements restricted to a smaller
range $m$ relative to $n$, and does not supply the present bound at
$m=n$. Weingartner \cite[Theorem 3 and Question 1]{We} studies the
version excluding $1$ and gives a lower estimate; it is not an upper
bound superseding \eqref{eq:chen}. The note \cite[Theorem 2.3, Corollary 2.6,
Proposition 2.7]{SM} formulates a general occupation-weight bound and a
finite acyclic flow linear program for weighted antichain problems and
leaves a $1+o(1)$ certificate for this problem open \cite[Question 5.1]{SM};
it supplies methodology, not a numerical bound of the present form.

Version 1 \cite{V1} proved
\begin{equation}\label{eq:v1}
 \rho(n)\le B_0(n):=1.0087578638+\frac{87.5686331624}{n}
 \qquad(n\ge1).
\end{equation}
The present version improves the strongest bound we located in a
literature search of 12 October 2026, in the sense of a smaller
asymptotic leading constant. That comparison includes our own
version 1 as well as Chen (1996). The search was bounded, not exhaustive:
it covered the problem page and discussion, arXiv, targeted searches,
and accessible citation and formalization records. A citation index
was rate-limited, another did not resolve Chen's 1996 record, and the
community discussion archive was only partly searchable. No priority
claim follows from this search. Section~\ref{sec:comparison} states
exactly where the displayed finite-$n$ bounds improve one another.

\section{Weights and their total mass}\label{sec:weights}
For $1\le a\le n$ let $M_a=\{ka:1\le k\le n/a\}$. If $A$ is
lcm-separated then its sets $M_a$ are pairwise disjoint: an element of
$M_a\cap M_b$ would be a multiple of $\lcm(a,b)$ not exceeding $n$.
Consequently, if $z_m\ge0$ and
\begin{equation}\label{eq:dual}
 \sum_{k\le n/a}z_{ka}\ge\frac1a\qquad(1\le a\le n),
\end{equation}
then
\[
 \sum_{a\in A}\frac1a
 \le\sum_{a\in A}\sum_{m\in M_a}z_m
 \le\sum_{m=1}^n z_m.
\]
The last inequality uses both disjointness and nonnegativity. This is
weak duality, and it includes the constraint for $a=1$.

Fix $L\in\{12,60,420\}$, $T=40$, and $Q=10^6$. Let $D$ be the
increasing list of divisors of $L$. The certificate provides integers
$C_{j,d}\ge0$ for $1\le j<T$ and $d\mid L$, and we define
\[
 c_{j,d}=C_{j,d}/Q\quad(1\le j<T),\qquad
 c_{0,d}=0,\quad c_{j,d}=0\quad(j\ge T).
\]
For positive integers $N$ and $g\mid L$ put
\begin{equation}\label{eq:F}
 F(N,g)=\sum_{k=1}^N\frac{c_{\lfloor N/k\rfloor,\,g\gcd(k,L/g)}}{k},
 \qquad F(0,g)=0.
\end{equation}
We shall certify $F(N,g)\ge1$ for every $N\ge1$ and $g\mid L$.
Set
\begin{equation}\label{eq:z}
 z_m=\frac{c_{\lfloor n/m\rfloor,\gcd(m,L)}}m\qquad(1\le m\le n).
\end{equation}
For $N=\lfloor n/a\rfloor$ and $g=\gcd(a,L)$, writing $a=gb$ and
$L=gh$ gives $\gcd(b,h)=1$, hence
\[
 \gcd(ka,L)=g\gcd(kb,h)=g\gcd(k,h),\qquad
 \left\lfloor\frac{n}{ak}\right\rfloor=\left\lfloor\frac Nk\right\rfloor.
\]
For the floor identity, write $n/a=N+\theta$, $0\le\theta<1$, and
$N=qk+r$, $0\le r<k$; then $0\le r+\theta<k$.
Thus the left side of \eqref{eq:dual} is $F(N,g)/a$, including at $a=1$.

For the mass estimate, consider the positive integers
$k_1<\cdots<k_s$ in $(X,Y]$ in one residue class modulo $L$, where
$0<X<Y$. The empty case contributes zero. Otherwise $1/k_1\le1/X$,
and for $i\ge2$ the spacing is $L$ and
\[
 \frac1{k_i}\le\frac1L\int_{k_{i-1}}^{k_i}\frac{dx}{x}.
\]
Summing these disjoint intervals inside $(X,Y]$ yields
\begin{equation}\label{eq:massap}
 \sum_{i=1}^s\frac1{k_i}\le\frac1L\log(Y/X)+\frac1X.
\end{equation}
The cell $\lfloor n/m\rfloor=j$ is $(n/(j+1),n/j]$. There are exactly
$\varphi(L/d)$ residues $r\in\{1,\ldots,L\}$ with $\gcd(r,L)=d$;
this follows by writing $r=du$ and counting the residues coprime to
$L/d$. Multiply \eqref{eq:massap} by the nonnegative cell coefficients
and sum to obtain
\begin{equation}\label{eq:mass}
 \sum_m z_m\le\tau+\frac Kn,
\end{equation}
where
\begin{equation}\label{eq:tauK}
 \tau=\frac1L\sum_{j=1}^{T-1}\sum_{d\mid L}
 \varphi(L/d)c_{j,d}\log\frac{j+1}{j},\qquad
 K=\sum_{j=1}^{T-1}\sum_{d\mid L}(j+1)\varphi(L/d)c_{j,d}.
\end{equation}
These formulas give an explicit error constant, with no asymptotic
remainder.

\section{The finite check}\label{sec:finite}
\begin{lemma}\label{lem:increment}
For every positive integer $N$ and every $g\mid L$,
\[
 F(N,g)-F(N-1,g)
 =\sum_{\substack{j\mid N\\j\le T}}
 \bigl(c_{j,e}-c_{j-1,e}\bigr)\frac jN,
 \qquad e=g\gcd(N/j,L/g).
\]
\end{lemma}
\begin{proof}
For $k\le N-1$, the floor $\lfloor N/k\rfloor$ increases precisely
when $k\mid N$, and then by one. The new term $k=N$ is covered by
$c_{1,g}-c_{0,g}$. Substitute $j=N/k$. Terms with $j>T$ vanish, but the
term $j=T$ must be retained because $c_{T-1,e}$ need not vanish.
At $N=1$ the identity reads $F(1,g)=c_{1,g}$.
\end{proof}

Put $S=2^{64}$. For every increment in Lemma~\ref{lem:increment}, the
program \texttt{check.c} adds
\[
 \left\lfloor\frac{(C_{j,e}-C_{j-1,e})jS}{N}\right\rfloor
 \quad\hbox{to }\mathrm{lo}_g,\qquad
 \left\lceil\frac{(C_{j,e}-C_{j-1,e})jS}{N}\right\rceil
 \quad\hbox{to }\mathrm{hi}_g.
\]
Starting at zero, addition of these directed bounds preserves
\begin{equation}\label{eq:enclosure}
 \mathrm{lo}_g\le QS F(N,g)\le\mathrm{hi}_g
\end{equation}
at every step. Signed integer division implements the floor and ceiling
also for negative numerators. The test $\mathrm{lo}_g\ge QS$ proves
feasibility; $\mathrm{hi}_g<QS$ is a violation, and
$\mathrm{lo}_g<QS\le\mathrm{hi}_g$ is ambiguous and does not certify it.

For these arrays, $0\le C_{j,d}\le1500031<2^{21}$ and $T=40<2^6$.
The numerator has absolute value below $2^{91}$. For $N\le10^7<2^{24}$,
partitioning $1,\ldots,N$ into $\{1\}$ and at most $24$ dyadic blocks
gives $\sum_{k\le N}1/k\le25$. At an intermediate update, each term
has either its old or its new nonnegative coefficient, so the exact
scaled sum is at most $25\cdot1500031 S$. There have been at most
$40N$ rounded increments. Therefore every accumulator has absolute value
less than
\[
 25\cdot1500031\cdot2^{64}+400000000<2^{90}.
\]
The signed 128-bit calculation cannot overflow. The ordinary integer
indices also fit on a platform with 64-bit \texttt{long}.

The exact checks gave zero violations and zero ambiguous cases for all
$g\mid L$. The original verification and two independent reviews checked
$1\le N\le10^7$, a total of $420000000$ pairs $(N,g)$ across the
three certificates. The independent Node program deposited with this
version was run for $1\le N\le100000$, checking $4200000$ pairs.
This shorter range already suffices with the tails below.
The smallest normalized lower margin for $L=12$ and $60$ is
$4/10^6$, at $N=1$; for $L=420$ the minimum occurs at $(N,g)=(2586,1)$
and its numerator in units of $QS$ is $234328714177376032570$.
All are strictly positive. The Node program additionally compared
\eqref{eq:enclosure} against direct sum enclosures at
$N=1,2,5,11,40,41,997,2586$ for every $g$, a total of $336$ checks.
These checks also prove the direct lower enclosures at those points
are at least $QS$.

\section{An explicit tail estimate}\label{sec:tail}
Let $\B_1(u)=\{u\}-1/2$ and $\B_2(u)=\{u\}^2-\{u\}+1/6$.
Then $|\B_1|\le1/2$, $|\B_2|\le1/6$, and $\B_2'=2\B_1$ away
from integers; $\B_2$ is continuous at integers. For fixed $g$ put
$e(r)=g\gcd(r,L/g)$ and define
\begin{align*}
 F_\infty(g)&=\frac1L\sum_{j=1}^{T-1}\sum_{r=1}^L
 c_{j,e(r)}\log\frac{j+1}{j},\\
 h_j(x)&=\sum_{r=1}^L(c_{j-1,e(r)}-c_{j,e(r)})
 \B_1\left(\frac{x-r}{L}\right),\\
 \Psi_g&=\sum_{j=1}^T j\sup_x|h_j(x)|,\qquad
 R_g=\sum_{j=1}^{T-1}\sum_{r=1}^L
 c_{j,e(r)}\frac{L(j+1)^2}{6}.
\end{align*}
In particular $F_\infty(1)=\tau$.

\begin{lemma}\label{lem:tail}
For every positive integer $N$ and every $g\mid L$,
\begin{equation}\label{eq:tail}
 F(N,g)\ge F_\infty(g)-\frac{\Psi_g}{N}-\frac{R_g}{N^2}.
\end{equation}
\end{lemma}
\begin{proof}
Fix $j<T$ and $r\in\{1,\ldots,L\}$, and put
$X=N/(j+1)$, $Y=N/j$, $a=(X-r)/L$, $b=(Y-r)/L$,
$f(u)=1/(Lu+r)$. Since $a>-r/L\ge-1$, the integers $i\in(a,b]$
correspond exactly to the positive integers $k=Li+r\in(X,Y]$.
The pole of $f$ lies strictly to the left of $a$.
Integration by parts on the intervals between consecutive integers,
using the downward jumps of size $1$ of $\B_1$, gives
\[
 \sum_{a<i\le b}f(i)=\int_a^b f(u)\,du
 +\B_1(a)f(a)-\B_1(b)f(b)+\int_a^b\B_1(u)f'(u)\,du.
\]
The convention $\B_1(\text{integer})=-1/2$ includes $b$ and excludes
$a$ if either is integral. Here the integral is $L^{-1}\log(Y/X)$,
$f(a)=1/X$, and $f(b)=1/Y$. A second integration by parts gives the
following bound for the final integral $E$:
\begin{align*}
 |E|&\le\frac{|f'(a)|+|f'(b)|}{12}
       +\frac1{12}\int_a^b|f''(u)|\,du\\
 &=\frac{|f'(a)|}{6}=\frac{L}{6X^2}.
\end{align*}
Indeed $f'(u)=-L/(Lu+r)^2$ and $f''(u)=2L^2/(Lu+r)^3>0$, so the
last integral is $|f'(a)|-|f'(b)|$.

Multiply by $c_{j,e(r)}\ge0$ and sum over cells and residues. The
main terms give $F_\infty(g)$ and the absolute remainders sum to at
most $R_g/N^2$. The boundary contribution is
\[
 \frac1N\sum_{j=1}^{T-1}\sum_{r=1}^L c_{j,e(r)}
 \left[(j+1)\B_1\left(\frac{N/(j+1)-r}{L}\right)
       -j\B_1\left(\frac{N/j-r}{L}\right)\right].
\]
Shifting the index in the positive part, using $c_{0,d}=c_{T,d}=0$,
turns this into $N^{-1}\sum_{j=1}^T jh_j(N/j)$. Its absolute value
is at most $\Psi_g/N$, proving \eqref{eq:tail}.
\end{proof}

\paragraph{Exact evaluation.}
Each $h_j$ is $L$-periodic and affine on $(x_0,x_0+1)$ for integral
$x_0$, with slope $L^{-1}\sum_r(c_{j-1,e(r)}-c_{j,e(r)})$.
Its supremum in absolute value is the maximum over its value at $x_0$
and its left limit at $x_0+1$, for $x_0=0,\ldots,L-1$.
Thus $\Psi_g$ and $R_g$ are rational and computed exactly.
The NumPy integer matrix calculation in \texttt{tail2.py} is safe:
its entries, including the slope correction, have absolute value at
most $(L^2+2L)1500031<2^{39}$, within signed 64-bit arithmetic.
The independent Node calculation instead uses a recurrence. With
$\delta_r=C_{j-1,e(r)}-C_{j,e(r)}$ and $H(x)=2LQh_j(x)$,
\[
 H(x+1)=H(x)+2\sum_r\delta_r-2L\delta_{x+1}.
\]
It checks both endpoints and periodic closure exactly with integers.

For logarithms, put $y=1/(2j+1)$ and
\[
 \ell_j=2\sum_{i=0}^{39}\frac{y^{2i+1}}{2i+1},\qquad
 u_j=\ell_j+\frac{2y^{81}}{81(1-y^2)}.
\]
Integrating the geometric series for $1/(1-y^2)$ shows
$\ell_j\le\log((j+1)/j)\le u_j$: every omitted positive term has
denominator at least $81$. For $j\ge1$ the width is at most
$1/(36\cdot3^{81})$. The sum of coefficients in any $F_\infty$ or
$\tau$ calculation is at most $39\cdot1.500031<59$, so the total
logarithm enclosure has width below $59/(36\cdot3^{81})$.
Write $F^-_\infty(g)$ for the expression using $\ell_j$, and $\tau^+$
for \eqref{eq:tauK} using $u_j$. All comparisons of these rational
numbers are exact.

\section{Certificate values and conclusion}\label{sec:values}
Table~\ref{tab:tail1} gives exact rational data at $g=1$; its column
$m$ satisfies $F^-_\infty(1)-1\ge m$.
\begin{table}[ht]
\centering\small
\begin{tabular}{r|c|c|c|r}
$L$ & $m$ & $\Psi_1$ & $R_1$ & $N_0$\\\hline
12 & $4479665243/10^{12}$ & $16566037/750000$ & $1188294863/250000$ & 5185\\
60 & $2711678024/10^{12}$ & $2254215137/60000000$ & $10493256427/100000$ & 16255\\
420 & $1890528110/10^{12}$ & $3518369257/52500000$ & $232154535249/50000$ & 70483
\end{tabular}
\caption{Exact tail constants and retained sufficient onsets.}\label{tab:tail1}
\end{table}
In this order the exact values of $mN_0^2-\Psi_1N_0-R_1$ are
\[
 \frac{138313568659001}{120000000000},\qquad
 \frac{12842678216509}{15000000000},\qquad
 \frac{7574273319816137}{300000000000},
\]
all positive. The corrected $L=420$ values are used throughout.
For every divisor $g>1$ simultaneously, Table~\ref{tab:tailg} gives
$F^-_\infty(g)-1\ge m_*$, $\Psi_g\le P_*$ and $R_g\le R_*$.
These bounds check the errors as well as the limiting margin.
\begin{table}[ht]
\centering\small
\begin{tabular}{r|c|r|r|c}
$L$ & $m_*$ & $P_*$ & $R_*$ & $10^{12}(m_*N_0^2-P_*N_0-R_*)$\\\hline
12 & $33442881402/10^{12}$ & 17 & 5473 & 805467948259683450\\
60 & $21435519688/10^{12}$ & 26 & 118639 & 5122531725449792200\\
420 & $15172074849/10^{12}$ & 49 & 5068429 & 66850545939558828361
\end{tabular}
\caption{Conservative exact bounds for all $g>1$.}\label{tab:tailg}
\end{table}
Because $\Psi_g,R_g\ge0$, the expression
$F^-_\infty(g)-1-\Psi_g/N-R_g/N^2$ is nondecreasing for $N>0$.
The positive comparisons at $N_0$ therefore prove $F(N,g)\ge1$ for
all $N\ge N_0$ and all $g$. Each $N_0<100000$, so this overlaps the
finite check of Section~\ref{sec:finite}. There is no uncovered positive
integer $N$.

The objective calculation gives the upward-rounded rational decimals
in Table~\ref{tab:objectives}; the exact $K$ values appear in the theorem.
\begin{table}[ht]
\centering
\begin{tabular}{r|c|c|c}
$L$ & $\tau\le\tau^+\le$ & exact $K$ (decimal) & headline bound\\\hline
420 & $1.001890528110611$ & $4162.909154$ & $1.00189053+4162.91/n$\\
60 & $1.002711678024969$ & $638.003685$ & $1.00271168+638.004/n$\\
12 & $1.004479665243758$ & $133.852416$ & $1.00447967+133.853/n$
\end{tabular}
\caption{Every displayed upper bound is rounded upward.}\label{tab:objectives}
\end{table}
The gaps between the headline leading constants and the displayed
bounds for $\tau^+$ are respectively
$0.000000001889389$, $0.000000001975031$, and $0.000000004756242$.
The headline $1/n$ coefficients exceed $K$ by $0.000846$, $0.000315$,
and $0.000584$. These are positive exact decimals. For each integer
$n\ge1$, \eqref{eq:z} now satisfies every constraint \eqref{eq:dual},
and \eqref{eq:mass} proves Theorem~\ref{thm:main}.

\section{Exact comparisons with earlier bounds}\label{sec:comparison}
\begin{remark}[the bound from version 1]
Version 1 proves $\rho(n)\le B_0(n)=1.0087578638+87.5686331624/n$
for every integer $n\ge1$. Denote the three displayed theorem
bounds by $B_{12},B_{60},B_{420}$. Table~\ref{tab:crossings} states all
pairwise comparisons of these four functions. In each row, the bound
in the second column is strictly smaller precisely for integers $n$
at least the final-column value; it is strictly larger at every smaller
positive integer. There is no integer equality.
\end{remark}
\begin{table}[ht]
\centering
\begin{tabular}{c|c|c|r}
old & new & exact real crossing & first integer with new $<$ old\\\hline
$B_0$ & $B_{12}$ & $77140611396/7130323$ & 10819\\
$B_0$ & $B_{60}$ & $2752176834188/30230919$ & 91039\\
$B_0$ & $B_{420}$ & $20376706834188/34336669$ & 593439\\
$B_{12}$ & $B_{60}$ & $50415100000/176799$ & 285155\\
$B_{12}$ & $B_{420}$ & $201452850000/129457$ & 1556138\\
$B_{60}$ & $B_{420}$ & $6408920000/1493$ & 4292646
\end{tabular}
\caption{Exact comparisons of the published decimal bounds.}\label{tab:crossings}
\end{table}
To justify the table, write $B_i(n)=a_i+b_i/n$. In every row
$a_{\rm old}>a_{\rm new}$ and $b_{\rm new}>b_{\rm old}$, and
\[
 B_{\rm new}(n)<B_{\rm old}(n)
 \quad\Longleftrightarrow\quad
 n>\frac{b_{\rm new}-b_{\rm old}}{a_{\rm old}-a_{\rm new}}.
\]
The displayed reduced rational crossings and their integer floors
were computed exactly. Thus the smallest of these four bounds is
$B_0$ for $1\le n\le10818$, $B_{12}$ for $10819\le n\le285154$,
$B_{60}$ for $285155\le n\le4292645$, and $B_{420}$ for
$n\ge4292646$. These comparisons concern the stated decimal
inequalities, not any unpublished sharpening of their constants.

For completeness Table~\ref{tab:chen} compares each with Chen's
explicit function $B_C$ on its domain $n\ge59$. Outside the intervals
listed, the bound in the first column is strictly larger than $B_C$;
there is no integer equality.
\begin{table}[ht]
\centering
\begin{tabular}{c|l}
bound & integers $n\ge59$ for which it is strictly below $B_C(n)$\\\hline
$B_0$ & $59\le n\le69$ or $n\ge8872$\\
$B_{12}$ & $59\le n\le64$ or $n\ge9537$\\
$B_{60}$ & $n=59$ or $n\ge43606$\\
$B_{420}$ & $n\ge274275$
\end{tabular}
\caption{Comparison with Chen (1996).}\label{tab:chen}
\end{table}
Indeed $n(n-58)(B_i(n)-B_C(n))$ is the concave quadratic
\[
 (a_i-1.0170166)n(n-58)+b_i(n-58)-14.21n.
\]
Exact opposite signs at the adjacent endpoints in the table locate
its two real roots; for $B_{420}$ the lower root lies in $(58,59)$.
A concave quadratic is positive between its roots and negative
outside. The accompanying audit records all endpoint signs as exact
fractions. This proves every stated comparison.

\begin{remark}[small values]
The exhaustive searches reported with version 1 found $\rho(5)=31/30$,
$\rho(11)=4699/4620$, and $\rho(n)=1$ for every other $n\le29$.
This remains an observation from that version, not a theorem or a new
computation in this version; its search programs are not deposited.
None of the present proof depends on it. The values at $5$ and $11$
also occur in the literature cited above.
\end{remark}

\section{Reproduction and verification status}\label{sec:repro}
The certificate files are \texttt{code/v2/cert\_L420.json},
\texttt{cert\_L60.json}, and \texttt{cert\_L12.json}. Their fields
\texttt{L}, \texttt{T0}, \texttt{Q}, \texttt{D}, \texttt{C} mean
$L$, $T=40$, $Q=10^6$, the ordered divisor list, and
\texttt{C[j-1][i]}$=C_{j,D_i}$ (zero-based column index $i$).
Their respective numbers of nonzero entries are $282$, $167$, and $99$.
All other entries are explicit zeros. The certificate arrays are
unchanged from those reviewed.

Run \texttt{sh code/v2/run.sh} from the bundle directory. It compiles
\texttt{check.c}, then invokes \texttt{verify.py} for each certificate
with $N_1=10^7$. Requirements are a C compiler supporting signed
\texttt{\_\_int128}, a platform with 64-bit \texttt{long}, Python 3,
and NumPy. No optimizer is required. The tail calculation uses exact
rational logarithm bounds and integer arrays. Its floating-point
fields and the C minimum-margin displays are supplementary diagnostic
images only. The final verified inequality uses upward decimal
rounding of the exact rational values of $\tau^+$ and $K$.
The expected-output notes distinguish predictions for that run from
outputs actually produced in this preparation environment.

An independent exact calculation is supplied in \texttt{audit.mjs}.
Run \texttt{node code/v2/audit.mjs}; its default finite limit is
$100000$. The saved \texttt{audit-output.json} records that run,
the exact tail values, directed objective decimals, pairwise bound
comparisons, and resource use. This is an independent implementation,
not an execution of the C/Python files. The preparation environment
had Node but no C compiler, Python, or TeX installation. A later
independent check compiled \texttt{check.c} and ran \texttt{verify.py}
on all three certificates up to $N=10^7$; its output equals the block in
\texttt{code/v2/expected-output.md}. The supplied \texttt{main-v1.pdf} is
the archived nine-page version 1, not this manuscript. The version 1
programs and certificate remain unchanged in \texttt{code/v1/}.

The certificate proof received two independent reviews with the
verdict PASS-WITH-REPAIRS and a separate exact verifier check.
The reviews requested documentary repairs: correcting the two
$L=420$ tail entries while retaining onset $70483$; restoring four C
and five Python string escapes damaged in source storage; and printing
certified bounds with upward rounding. The two independently repaired
source copies agree byte for byte. The deposited C and tail sources
are those copies. The Python wrapper additionally has the output-only
change recorded in \texttt{printing.patch}, and an internal provenance
tag was removed from a comment. The feasibility tests, certificate
arrays, tail inequalities, and verification decision are unchanged.
The original verifier check used the same C source with an alternative
C compiler; the independent reviews also replayed the full finite
range in exact integer implementations. The present Node run provides
a further, shorter exact check with complete tail overlap. No formal
proof-assistant verification is claimed.

The arrays were found by floating-point linear-programming searches
followed by rational rounding and feasibility checks. These searches
and any floating-point sanity checks are supplementary; the proof is
the exact certificate verification described above. The periodic-dual
method suggests limit linear programs as a possible subject for
further study.

\begin{thebibliography}{9}
\bibitem{EP542}
T. F. Bloom, \emph{Erd\H{o}s Problems}, Problem 542,
\url{https://www.erdosproblems.com/542}, and discussion
\url{https://www.erdosproblems.com/forum/thread/542},
accessed 12 October 2026.
\bibitem{Ch96}
Y.-G. Chen, \emph{On a problem of P. Erd\H{o}s},
Acta Scientiarum Mathematicarum (Szeged) \textbf{62} (1996), 101--114.
\url{https://acta.hu/download.phtml?id=2350}.
\bibitem{Ch97}
Y.-G. Chen, \emph{On a paper of Schinzel and Szekeres},
Acta Scientiarum Mathematicarum (Szeged) \textbf{63} (1997), 369--379.
\url{https://acta.hu/download.phtml?id=2393}.
\bibitem{We}
A. Weingartner, \emph{The Schinzel--Szekeres function},
Research in Number Theory \textbf{11}, 63 (2025);
arXiv:2310.13038, version 2, 13 June 2025.
\url{https://arxiv.org/abs/2310.13038}.
\bibitem{SM}
\emph{Sub-Markov chain certificates for weighted antichain bounds},
research note, 15 April 2026,
\url{https://www.ulam.ai/research/markov-chain-erdos.pdf}.

\bibitem{V1}
H. Chen, \emph{An explicit upper bound for reciprocal sums of sets
with large pairwise least common multiples}, version 1,
12 October 2026, Zenodo.
\url{https://doi.org/10.5281/zenodo.23311400}.
\end{thebibliography}
\end{document}
